\documentclass{article}
\usepackage[UTF8]{ctex}
\usepackage{amsmath}
\usepackage[numbers,sort&compress]{natbib}
\usepackage{graphicx}

\title{政策意义侵蚀场景测试报告}
\author{郑子薇}
\date{19th March 2026}

\begin{document}

\maketitle

\section{场景设定背景}

“政策意义侵蚀”场景旨在模拟一项政策命令在多层级组织中自上而下传递时，如何在解释、转述、筛选与执行的连续过程中发生语义漂移、执行收缩与选择性保留。该场景关注的并非政策文本本身是否合理，而是政策在组织链条中如何由“正式文本”逐步转化为“可执行文本”，并最终在基层形成与原始政策存在差异的行动版本。\cite{PressmanWildavsky1984}\cite{Lipsky2010}

这一问题对应公共政策执行研究中的经典命题，即政策并不是在颁布之时就自动完成，而是在不同层级、不同岗位、不同资源约束与责任压力下，被不断再解释、再编码与再实施。\cite{PressmanWildavsky1984}\cite{SabatierMazmanian1980} 政策传递链条越长、参与者裁量空间越大、考核压力与资源缺口越显著，原始政策含义越可能被压缩、替换、弱化，甚至被截留。\cite{HillHupe2002}\cite{Lipsky2010}

与仅仅把组织理解为“自上而下的整层广播系统”不同，本场景进一步将政策执行过程放回到具体的组织连接结构之中。现实中的政策往往并不是由某一层一次性、均匀地传递给下一整层，而是沿着部门、条线、分管关系和工作接口所构成的局部网络继续扩散。某些节点拥有明确上游来源，某些节点则在当前政策回路中并未被激活；某些节点仅受单一上级影响，而某些节点则需要同时吸收多个上级版本并加以综合。由此，政策“意义侵蚀”不仅是层级效应，也是网络结构效应。\cite{HillHupe2002}\cite{Weick1995}

从现实政策情境看，这种“意义侵蚀”现象常见于层级治理、专项整治、绩效考核、风险排查与合规执行等任务型治理情境中。高层发布的政策往往包含目标、范围、标准、资源、报送、责任与不可变更条款；但中层出于协调成本、问责压力与绩效约束，会倾向于将其重组为部门化任务；基层则进一步将其压缩为最低可执行动作。在这一过程中，政策并非简单地被“误传”，而是被组织成员根据本层利益结构、执行可行性和所在网络位置进行“有方向的重构”。\cite{Lipsky2010}\cite{Weick1995}

当前版本的场景设定已不再局限于“单向层级失真”这一最小模型，而是试图将政策意义变化放回到更接近真实组织政治过程的循环之中。也就是说，政策意义侵蚀不仅可能发生在 top $\rightarrow$ mid $\rightarrow$ low 的正式传递阶段，也可能在政策进入执行阶段之后，继续通过上行反馈、同层协商、非正式传播、越级投诉、转办下级与再调整而被重新塑造。这样，实验所关注的对象便不再只是“政策如何往下传”，而是“政策如何在组织中被持续加工、再编码与再政治化”。

\section{场景设计的理论依据}

本场景的设计主要建立在四条理论脉络之上。

第一，政策执行研究中的“执行落差”命题强调，政策目标与执行结果之间存在系统性偏差。Pressman 与 Wildavsky 指出，政策落实并非单次命令传递，而是一个跨组织、跨层级、跨环节的复杂协同过程；执行链条上的每一次转换，都会引入新的不确定性与偏差。\cite{PressmanWildavsky1984} 因此，本场景将政策传递设计为严格的层级链条，而非所有主体同时接收同一原文，以便观察链式传播中的累积失真。

第二，街头官僚理论强调，一线执行者并不是被动执行者，而是在资源不足、任务模糊与责任压力下，依靠裁量权将抽象政策转化为具体实践的人。Lipsky 的研究指出，基层行动者在高负荷与有限资源条件下会发展出简化、筛选、拖延与例行化的策略，从而使“政策文本中的国家”转化为“实践中的国家”。\cite{Lipsky2010} 本场景的失真模式以及后续反馈模式，正是将这种裁量性嵌入到下传、截留、反馈与再解释机制中。

第三，公共政策实施中的治理与中介层级研究指出，中层组织并非纯粹的传声筒。Hill 与 Hupe 认为，政策执行是治理结构中的多重互动过程，中层承担任务拆解、口径协调、资源再分配与责任转译等职能。\cite{HillHupe2002} 因此，本场景中层并不被建模为机械复制者，而被赋予较高的重述、筛选、压缩与反馈处置倾向，以反映其在实际治理中的“翻译者”与“过滤者”角色。

第四，组织中的意义建构理论说明，组织成员不会机械接受信息，而是依据其所处位置、职责、风险感知与既有认知框架，对外部命令进行解释。Weick 强调，组织中的行动者总是在不确定环境中通过意义建构形成行动方向。\cite{Weick1995} 本场景中，高层、中层与基层拥有不同的语义优先级、重写偏好以及对风险、责任与可执行性的不同敏感度，正是为了把这种分层意义建构机制操作化。

在此基础上，本场景进一步引入“层级约束下的网络可达性”与“级联结束后的后续互动”两类机制。前者说明政策不是向下一整层平均扩散，而是在层级顺序约束下，只沿着实际存在的组织连接继续传播；后者说明政策传递并不在基层接收后自动结束，而是可能通过执行困难反馈、越级投诉、同层协商、上级再发布与定向私有重启等方式，进入新一轮意义重构。这样一来，政策执行研究中的“下传链条”与治理研究中的“反馈回路”便被整合进同一实验框架中。\cite{PressmanWildavsky1984}\cite{HillHupe2002}\cite{Weick1995}

综上，本场景并非简单模拟“信息噪声”，而是将政策执行研究中的层级传递、街头裁量、组织翻译、意义建构、连接结构选择效应以及反馈回路整合为一个可计算、可比较、可复现实验场景。\cite{PressmanWildavsky1984}\cite{Lipsky2010}\cite{HillHupe2002}\cite{Weick1995}

\section{场景设计}

\subsection{场景的具体结构与运行逻辑}

在当前实现框架中，场景采用“层级顺序 + 网络路由 + 后续互动”的三重机制，如图~\ref{fig:kuangjia}所示。系统默认的层级顺序为 top $\rightarrow$ mid $\rightarrow$ low，但实际的政策传播并不是在某一层完成后自动广播到下一整层，而是在层级顺序允许的前提下，只沿着研究者预先设定的组织连接继续传播。由此，场景的基本单位不再是“整层同时接收同一政策版本”，而是“在层级边界内沿可达分支逐级传递”。

\begin{figure}[htbp]
    \centering
    \includegraphics[width=\linewidth]{政策意义侵蚀系统框架.png}
    \caption{场景框架图}
    \label{fig:kuangjia}
\end{figure}

\begin{figure}
    \centering
    \includegraphics[width=1\linewidth]{层级设置.png}
    \caption{政策层级设置}
    \label{fig:tiersetting}
\end{figure}

\begin{figure}
    \centering
    \includegraphics[width=1\linewidth]{层级预览.png}
    \caption{层级设置预览}
    \label{fig:tierview}
\end{figure}

若记组织网络为有向图
\[
G=(V,E),
\]
其中 $V$ 表示行动者集合，$E$ 表示从发送者到接收者的可传播连接；再记层级映射为
\[
\ell:V\rightarrow\{\texttt{top},\texttt{mid},\texttt{low}\},
\]
则某一行动者 $u$ 的合法下游集合可表示为
\[
D(u)=\{\,v\in V\mid (u,v)\in E,\ \ell(v)=\text{next}(\ell(u))\,\}.
\]
也就是说，网络边决定“谁可以传给谁”，层级顺序决定“现在允许传到哪一层”。只有同时满足这两项条件，目标节点才会在下一阶段成为可激活对象。

在正式级联阶段，政策作为最新公告或广播进入系统，并自上而下沿层级边界与网络可达路径逐步扩散。若某节点同时具有多个上游来源，则该节点接收到的并不是一个天然统一的政策文本，而是多个版本的上游输入，随后由该节点根据自身角色职责和解释偏好形成统一的本层表达。由此，网络汇合点不再只是一个信息通过点，而成为政策意义再次重组的关键位置。

与早期的“固定回合数前进”不同，当前实现已经把“当前活动级联层必须跑空”作为运行约束之一。也就是说，只要当前 cascade tier 中仍有被激活但尚未完成处理的节点，系统就会继续推进，直到同层目标全部完成后才允许停止。这一调整的意义在于，实验中不再会因为顺序器的有限预算、空转回合或被跳过的非活动节点而让同层某个分支在日志上“凭空消失”。因此，当前的层级传播不再只是静态的顺序定义，而是包含一个明确的运行时完结条件：同层目标必须被完整清空。

需要特别说明的是，正式级联结束之后，系统并不会机械地不断重复下一轮级联。若在级联完成后没有新的公共政策广播进入系统，则场景进入后续互动阶段。此时，各节点不再默认触发新的自上而下政策传递，而是围绕既有政策文本和执行后果形成上行反馈、升级反馈、越级投诉、同层协商、转办下级与非正式传播等特殊事件。换言之，正式政策传递结束后，政策意义的组织性变化并未结束，而是转入执行过程中的反馈性再建构。

与此同时，若系统在后续互动阶段再次接收到一条新的政策广播，则这一新的政策输入会终止当前后续互动回路，并重置场景的活动状态，使系统重新以该新政策为最新政策文本，进入新的层级级联轮次。当前版本中，这种“重启”又进一步细分为两类。其一，若外部系统重新广播一条公共或定向私有政策，则系统直接以该政策重置活动状态；其二，若顶层在后续互动阶段通过 \texttt{announce\_policy\_adjustment} 发布新的政策调整，则系统不再只重新激活该发布者自身下游的单一分支，而是直接重新激活整个下一层，从而把新的顶层政策视为一项针对该层级全体对象的重新部署。由此，后续互动模式并不是一个与正式级联并行存在的永久状态，而是夹在两轮正式政策广播之间的过渡性执行阶段；一旦新的广播政策到来，场景的主导逻辑便重新回到政策自上而下的层级传递。

对于定向私有政策广播，当前版本还增加了“按接收者逐条可见”的事件记录机制。也就是说，当一条定向私有级联输入同时发给多个顶层节点时，系统会分别生成针对每个接收者的环境可见事件，而不是只保留第一名接收者的可见记录。这一修改使实验日志与实际激活状态一致，避免出现“系统显示私有广播发给多人，但环境事件只写其中一人可见”的观察偏差。

\subsection{网络连接与层级耦合的实验含义}

网络连接机制的引入，使本场景从“单链式层级传递”扩展为“分支化层级传递”；后续互动机制的引入，则进一步使场景从“单向命令下传”扩展为“多方向、非同步、可回流的组织互动”。如果说传统的层级级联主要回答“政策在越往下传时是否越容易失真”，那么当前版本还可以进一步回答“在不同部门、不同分管线路和不同连接结构中，失真是否会以不同方式发生，并在执行阶段如何被重新放大、遮蔽或修正”。

例如，在一个三层组织中，研究者可以将顶层中的两个节点分别设定为经济分支和教育分支的起点，再通过网络边将其连接到各自的中层和基层子树。此时，实验中的传播过程不再是
\[
\texttt{top}\rightarrow \texttt{mid}\rightarrow \texttt{low}
\]
意义上的层级平均扩散，而更接近
\[
\texttt{top}_1\rightarrow \texttt{mid}_1\rightarrow \texttt{low}_{1a},\texttt{low}_{1b},
\qquad
\texttt{top}_2\rightarrow \texttt{mid}_2\rightarrow \texttt{low}_{2a},\texttt{low}_{2b}.
\]
在这种结构下，不同部门所形成的政策文本，不仅会因为角色属性不同而分化，也会因为上游来源不同而形成分支化演化路径。于是，研究者可以比较的就不只是“高层与基层的差异”，而是“不同政策支路在相同层级上的差异”。

若进一步考虑后续互动，则同一分支内部还可能出现基层向中层的反馈、基层向高层的升级或越级投诉、以及同层行动者之间的非正式协商。由此，政策意义侵蚀可以被分解为四类来源：其一，层级位置带来的结构性压力；其二，角色属性带来的解释偏好；其三，组织连接带来的路径依赖；其四，后续互动带来的循环性重构。

\subsection{两种实验模式及其区别}

本场景包含两种核心模式：严格级联模式（\texttt{strict\_cascade}）（图~\ref{fig:strict}）与失真级联模式（\texttt{distortion\_cascade}）（图~\ref{fig:distortion}）。
\begin{figure}
    \centering
    \includegraphics[width=1\linewidth]{严格.png}
    \caption{严格级联模式}
    \label{fig:strict}
\end{figure}
\begin{figure}
    \centering
    \includegraphics[width=1\linewidth]{失真.png}
    \caption{失真级联模式}
    \label{fig:distortion}
\end{figure}
严格级联模式用于构造政策忠实传递的基准条件。在该模式下，系统要求上层向下层传递时尽量保留原始政策内容，并且通过规则校验保证政策原文、关键硬约束或前段主体内容不会被任意删改。行动者仍可补充本层职责相关的执行细节，但这种补充必须服从“原文优先”的约束。其研究意义在于提供一个低裁量、高保真的对照组，用于观察即便在技术上强约束保真的条件下，不同层级与不同分支是否仍会因为职责差异而形成解释偏移。\cite{PressmanWildavsky1984}

失真级联模式则用于构造更贴近真实组织执行过程的条件。在该模式下，行动者不仅可以忠实转述，也可以基于本层利益、风险压力与执行能力，对政策进行选择性强调、弱化、压缩、重写，甚至直接截留。此时系统不再强制要求保留完整原文，而是引入若干参数控制“改写程度”“冲突敏感度”与“截留倾向”。因此，失真模式更适合研究以下问题：其一，层级越低是否越倾向于把政策压缩为最低执行口径；其二，考核压力、资源缺口与执行负担如何推动政策含义向“可操作但偏离原意”的方向演化；其三，不同网络分支是否会产生不同的政策压缩路径；其四，在哪些参数组合下，政策会在特定分支的中层或基层被系统性截留。\cite{Lipsky2010}\cite{HillHupe2002}

在当前版本中，两种模式的差异还进一步延伸到了后续互动阶段。严格级联模式下，后续互动更多表现为口径澄清、信息补充与有限反馈，系统更倾向于维持原始政策的连续性；失真级联模式下，后续互动则更容易转化为资源抱怨、执行困难上报、同层规避性协商、越级投诉与政策再次压缩。换言之，模式差异不仅解决“如何改写”，也影响“级联完成后如何继续互动”。

\subsection{失真模式的参数化设计及其计算逻辑}

失真模式下，场景使用三个核心参数控制政策意义侵蚀的发生概率与侵蚀幅度，即失真强度（\texttt{distortion\_strength}）、冲突敏感度（\texttt{conflict\_sensitivity}）和阻断概率（\texttt{block\_probability}）（图\ref{fig:distortionpara}）。在当前实现中，这三个参数的默认值分别为
\[
s = 0.60,\qquad \lambda = 0.50,\qquad b = 0.25.
\]
\begin{figure}
    \centering
    \includegraphics[width=1\linewidth]{失真参数.png}
    \caption{失真模式参数设置}
    \label{fig:distortionpara}
\end{figure}
其中，失真强度 $s$ 控制文本重写的幅度；冲突敏感度 $\lambda$ 控制行动者对政策—角色冲突的放大程度；阻断概率 $b$ 则给出“不继续下传、直接截留”的基础倾向。

系统首先根据层级位置给出一个基础压力值：
\[
\tau_{\text{top}} = 0.18,\qquad
\tau_{\text{mid}} = 0.42,\qquad
\tau_{\text{low}} = 0.68.
\]
这意味着，即使不考虑具体政策内容，仅从组织位置出发，基层也被设定为承受更高的执行压力与冲突风险。\cite{Lipsky2010}

随后，系统从政策文本中提取若干语义信号，例如执行负担、报送要求、问责强度、资源支持与硬约束条款，并与角色画像中的负担感、资源依赖、自主性偏好、控制取向和稳定性偏好进行耦合，形成语义冲突项 $C_i$。在形式化表达上，这一过程可写为
\[
C_i=\operatorname{clip}\Big(
w_1\cdot \text{burden}
+w_2\cdot \text{resource\_gap}
+w_3\cdot \text{accountability}
+w_4\cdot \text{report}
-\delta_1\cdot \text{control}
-\delta_2\cdot \text{stability}
+\eta_i
\Big),
\]
其中 $\eta_i$ 表示层级修正项：高层获得负修正，基层获得正修正。\cite{HillHupe2002}\cite{Lipsky2010}

在此基础上，系统计算冲突压力：
\[
P_i=\operatorname{clip}\Big(\tau_i(1-\lambda)+C_i\lambda\Big).
\]
当 $\lambda$ 较低时，压力主要由层级基础位置决定；当 $\lambda$ 较高时，压力更多由具体政策内容与角色特征之间的语义冲突决定。

接着，系统构造一个阻断激活系数：
\[
A=\operatorname{clip}\big(0.10+0.55s+0.35\lambda\big).
\]

最终，系统计算某一行动者在当前轮次的阻断倾向：
\[
B_i=\operatorname{clip}\big(b+0.55P_iA+0.05\epsilon_i\big),
\]
其中 $\epsilon_i\in[0,1]$ 为一个由行动者身份与当前政策状态共同决定的确定性微扰项。

当
\[
B_i \ge 0.5
\]
时，系统判定该行动者在本轮选择截留政策，而不是继续向下传递。\cite{Lipsky2010}

需要强调的是，在当前版本中，参数控制的是“到达某个节点后的改写与截留倾向”，而网络连接控制的是“该节点是否能够被到达”；而在后续互动阶段，持续制度条件则进一步影响反馈、协商、升级与回避的倾向。因此，一个节点不响应，可能既意味着它在正式级联中从未被上游路径命中，也可能意味着它在后续互动阶段因高压力或低收益判断而选择沉默。将这两种机制区分开来，是当前场景的重要方法论收益。

\subsection{失真前后文本的对照显示机制}

在当前版本中，失真级联模式不仅关注节点最终生成的下传文本，还提供了失真前后政策文本的对照显示机制（图~\ref{fig:distortioncomparsion}）。该机制并非简单地把两段文本并排呈现，而是将“当前节点收到的上游政策版本”与“该节点最终真正发给下一级的版本”以左右双栏形式并置显示，并进一步对新增、删除与改写部分进行逐行和行内高亮标记。由此，研究者不仅能够看到最终结果是什么，还能够直接观察政策意义究竟是在何处发生了压缩、替换、删减或重写。

更具体地说，在该对照机制中，左侧显示的是当前层级收到的上级政策版本，右侧显示的是该节点经过处理后真正向下级发送的内容。若系统同时保留了该节点的原始生成草稿，则还可以在附加区域中显示该草稿文本。这样，研究者便能够同时比较三个层次：其一，上游输入版本；其二，agent 原始草稿版本；其三，最终实际下传版本。于是，失真不再只是“输入与输出不同”这一抽象判断，而是被进一步分解为“agent 自发改写了什么”与“系统最终保留了什么”这两个可区分的过程。

从实验解释的角度看，这一对照显示机制至少具有三方面意义。第一，它显著提高了政策意义侵蚀的可解释性，使研究者能够直接定位具体失真发生在哪些条款、哪些表述和哪些责任链条上。第二，它有助于区分“表面上仍然可执行”的文本与“形成过程中已经明显偏离原文”的文本，从而避免只根据最终文本外观低估失真程度。第三，它使研究者能够将失真进一步分类为新增、删除与改写三种基本操作，从而支持更细粒度的实验统计与结果比较。

因此，在当前场景中，研究者不应只记录节点最终发出的政策文本，还应结合失真前后的对照结果，考察各节点在政策目标、适用范围、执行标准、资源支持、报告要求、责任分工与不可改写条款等方面究竟保留了什么、弱化了什么、替换了什么。这样，文本失真就不再只是一个笼统的结果判断，而成为一个可以逐轮、逐节点、逐条款观察和分析的可解释过程。
\begin{figure}
    \centering
    \includegraphics[width=1\linewidth]{失真对比.png}
    \caption{失真模式日志对比}
    \label{fig:distortioncomparsion}
\end{figure}

\subsection{后续互动模式的事件结构}

在当前设计中，正式级联结束后的场景并非处于空转状态，而是进入一个以私有互动为基本单位的后续模式。该模式的核心特征是：在没有新的广播政策进入系统时，节点不再自动触发新一轮层级下传，而是围绕上一轮政策及其执行后果形成各种特殊事件。

第一类是正常反馈。基层或中层可以向直属上级发起关于执行成本过高、资源不足、目标冲突、群体反弹、无法按时完成等问题的上行报告。此类反馈进入上级的私有输入，而不被视为全局公告。其机制意义在于将“执行困难”从一段文本内容，转化为具有明确发送者、接收者与待处理状态的组织事件。

第二类是升级反馈。若直属上级未处理、延迟处理、压制或歪曲前述反馈，则发起者可以在后续轮次中继续向更高层升级报告。升级反馈并不保证一定改变政策，而是将未经中层过滤的版本送入更高层的判断空间，从而破坏中层对信息的垄断。

第三类是越级投诉或越级告状。与一般升级反馈不同，越级投诉是专门指向跨层级直接报告的特殊事件。其触发往往与明显违规风险、不可改写条款被删除、资源条件与执行要求严重失配、连续多轮反馈无回应等情形有关。其制度含义不在于“多发一条消息”，而在于暂时中断原有的层级过滤秩序，使高层得以接触来自基层的原始风险版本。

第四类是同层协商与非正式传播。若两个行动者处于同一层级且在当前组织连接或非正式连接中彼此可达，则在无新广播政策时，其中一方可以向另一方发起私下协商、经验交换、口径核对或非正式提醒。这一机制使场景中的政策解释不再完全由正式层级结构决定，而开始受到同层关系网络与私下互动的影响。

第五类是下级转办与定向续办。中层或高层在私有线程中并不只面对“向上答复”这一单一路径，而可以把某些问题私下通知下级，要求其补充数据、执行排查或完善解释口径。这样，后续互动不只是“反馈向上”，也可能表现为“问题沿私有线程被重新压回下级执行”。

第六类是等待性停顿。当前版本已经把“当前已没有任何动作倾向，建议注入新的环境事件或发布新的政策”设定为一个明确的无动作句，并允许系统在后续阶段把已明确声明无动作倾向、同时又没有新的私有线程输入的节点置于等待状态。其含义是：后续互动不是无穷无尽的强制对话，而是在保持可继续反馈空间的同时，允许某些节点进入“等待下一次公共事件或新政策”的静默阶段。

因此，当前版本中的后续互动不再只是一个模糊的“补充讨论阶段”，而是已经具有明确动作类型、线程状态、接收对象、转办方向与静默条件的私有组织互动结构。

\subsection{持续制度条件与后续干预方式}

在当前版本中，环境变量不再被统一视为一次性事件，而是被区分为短期事件与持续制度条件两类（图~\ref{fig:ganyu}）。
\begin{figure}
    \centering
    \includegraphics[width=1\linewidth]{对照实验干预.png}
    \caption{对照实验干预}
    \label{fig:ganyu}
\end{figure}
短期事件主要指突发舆情、上级督查、媒体曝光、群体投诉等瞬时性冲击。此类事件仍可通过广播注入系统，并在当轮或近几轮中改变行动者判断。

持续制度条件则不同。资源紧张、监察常态化、舆论高压、考核周期临近、长期人员不足、地方保护压力等，更适合被建模为持续存在的结构性背景，而不是一条“新的通知”。因此，当前版本允许研究者将这些条件直接写入场景状态，使其在多轮中持续参与冲突压力计算、截留倾向估计、上行反馈概率判断与后续互动提示生成。

这一设计带来一个重要方法论收益：研究者可以在不触发新一轮层级级联的情况下，只干预后续互动阶段。例如，在正式级联完成后，通过提高资源短缺程度或舆论压力水平，观察基层是否更倾向于发起执行困难反馈，中层是否更倾向于压缩口径，高层是否更可能选择忽略、转办或再次澄清。由此，环境干预与层级级联被区分为两个不同维度的实验操作。

当前版本还进一步强化了“冲击—线程—再发布”之间的联动关系。也就是说，突发舆情、抗议游行、媒体曝光等环境事件不仅更新短期判断，还会影响后续线程提示中的议题焦点、上报倾向与再调整压力；在达到一定阈值时，顶层又可能基于这些后续信号重新发布政策调整，从而让持续制度条件和短期冲击真正进入下一轮级联结构。

\section{真实实验测试}

\subsection{实验一：默认实验测试-qwen3:0.6b}

\subsection{实验二：aaa}
\subsection{实验三：bbb}
\subsection{实验四：ccc}

\section{目前版本局限}

尽管当前场景已经能够较好地模拟政策在层级结构中的分支化传递、意义侵蚀与执行阶段再解释，但从真实社会模拟的角度看，现有实现仍然存在若干重要局限。

第一，当前场景虽然已经引入了后续线程、定向私有广播、同层协商、升级反馈与再发布机制，但总体骨架仍然以较强的层级顺序逻辑为核心。真实组织中更常见的跨层并行推进、非同步响应与部分节点反复介入，在现有框架中仍然被明显压缩。当前代码版本下，场景的核心状态变量仍然围绕 \texttt{current\_tier\_idx}、\texttt{tier\_seen}、\texttt{active\_tier\_targets} 和 \texttt{latest\_policy} 展开，并通过当前活动层级控制谁可以行动，这说明模型的骨架仍然是一个较强的顺序级联框架，而不是一个真正意义上的异步多线程组织过程。

第二，当前实现中的行动空间已明显比早期版本丰富，因而“只有 \texttt{send\_message} 与 \texttt{yield} 两类动作”这一局限已经不再成立。当前版本中，场景已经把 \texttt{report\_upward}、\texttt{escalate\_complaint}、\texttt{consult\_peer}、\texttt{notify\_subordinate} 与 \texttt{announce\_policy\_adjustment} 等后续互动动作纳入正式运行逻辑。然而，这并不意味着行动空间已经足够贴近真实组织实践。与现实中的请示、会签、联合报送、会议纪要确认、口头承诺、策略性拖延、形式合规与实质规避相比，当前动作集合仍然偏少，很多组织性互动还未被拆解为更细粒度、可验证的制度化动作。

第三，当前场景对组织连接的刻画虽然已经超出“整层平均广播”，并区分了正式传播路径、同层非正式协商与私有线程可见关系，但对多重网络的并行共存仍然刻画不足。现实组织中往往同时存在正式指挥链、业务协同链、私人关系链、监督举报链和跨部门接口链；当前模型虽然已经把部分关系拆开，但距离把这些不同网络分别作为独立约束结构纳入计算，仍有明显距离。未来若要提升社会模拟真实性，更合理的方向是把正式传播边、非正式传播边、汇报边、协商边和监督边分开建模，而不是继续让多类互动共享较少数的连接机制。

第四，当前模型对时间结构的刻画仍然较弱。现有逻辑更多是“轮次推进”而非“时序约束”。这意味着现实政策执行中非常重要的期限、窗口期、延迟反馈、积压效应、突发事件插队与考核节点冲击，目前仍然主要表现为提示语义上的压力，而不是严格进入事件调度、等待时间和决策时点的计算过程。若未来要加强真实感，可以考虑把时限、超期惩罚、周期性检查和批次汇总机制纳入状态转移中，使“何时说”“多久不回应”“在什么时间点升级”为可观测变量。

第五，当前场景对政策内容本身的结构化处理仍有继续深化的空间。虽然当前设计已将政策拆分为目标、范围、标准、资源、报送、责任与不可改写条款等子条款，但真实政策执行中的争议往往不仅来自条款是否保留，还来自条款之间的优先级冲突、执行条件依赖和跨部门责任模糊。也就是说，政策并不只是若干平行语义段落的集合，还可能是一个包含前提、例外、触发条件与优先级排序的复杂规则系统。未来若能进一步把条款关系图纳入模型，场景对“表面未失真但实质已偏移”的识别能力会显著提高。

第六，当前场景对行动者认知与学习机制的处理仍然偏静态。现实中的组织成员并不会在每一轮都以相同偏好处理政策，相反，他们会根据前几轮的问责经验、协商结果、上级反应、社会舆论与个人得失不断调整自己的策略。但在当前实现中，角色画像更多仍是相对固定的属性集合，而非可随互动历史演化的策略系统。未来若希望增强社会模拟真实性，可以引入记忆累积、信誉变化、风险预期更新和策略学习机制，使行动者真正具备“被场景塑造”的动态性。

第七，当前场景对“解释权竞争”的实现仍然有待从概念层提升为可验证的计算结构。真实组织中的政策分支并不是简单并存，而往往围绕“谁更代表领导意图”“谁更安全”“谁更可执行”展开竞争。现阶段这一部分虽然已经在场景中通过多上游版本、同层协商和上游汇合点综合机制被部分引入，但若要在实验上形成更强的解释力，仍需要把分支竞争显式转化为可记录、可比较、可回溯的合法性评分、采纳概率或再澄清触发条件。

第八，当前场景的后续互动干预虽然提升了实验灵活性，但也带来了一个新的方法论问题，即实验设计者对场景后续过程的影响力仍然较强。若研究者可以直接注入私有线程、持续条件或后续互动种子，那么实验结果的一部分可能来自研究者的人工设定，而不是组织成员在既有制度环境中自然涌现的行为。未来若要增强模型的外部效度，需要在“可控干预”与“自然涌现”之间保持更清晰的边界，例如区分基准实验、干预实验与探索实验三类模式。

总体而言，当前“政策意义侵蚀”场景已经能够较好地模拟政策在层级组织中的链式传播、分支化失真、后续线程互动与再调整重启，但若要从“可运行的政策级联模型”进一步提升为“更高保真的真实社会模拟平台”，仍需在异步互动、多网络建模、动态学习、时间约束、解释竞争与制度复杂性方面继续推进。

\bibliographystyle{unsrtnat}
\bibliography{refs}

\end{document}
